\documentclass[11pt]{amsart}
\usepackage[margin=1in]{geometry}
\usepackage{amsmath,amssymb,amsthm}
\usepackage{enumitem}
\usepackage{booktabs}
\usepackage[hidelinks]{hyperref}
\setlist{nosep,leftmargin=*}

\newtheorem{theorem}{Theorem}
\theoremstyle{remark}
\newtheorem*{remark}{Remark}

\newcommand{\tn}[1]{[#1]}
\newcommand{\qq}{\mathfrak{q}}
\newcommand{\st}{\star}

\title[DS$(r,1,1)$ for all $r\ge 2$]{DS$(r,1,1)$ for all $r \ge 2$: a direct proof via Thm 3.12 and Sub-Lemma Z}
\author{Rick}
\date{2026-09-26}

\begin{document}

\begin{abstract}
We prove the dominance-support statement DS$(r,1,1)$ for every $r\ge 2$, in its sharpened triangular form and in the stable limit. We give an explicit closed form for $C_r = e_1\st e_1\st e_r$ in the $e$-basis. The proof uses only $e_1(Y)$: Hikita's Thm~3.12 in operator form, together with Sub-Lemma~Z. It does not need $W_r$. This is the first length-3 instance of DS that is a theorem.
\end{abstract}

\maketitle

\begin{center}
\begin{tabular}{ll}
\textbf{Author:} & Rick\\
\textbf{Date:} & 2026-09-26\\
\textbf{Recipient:} & Clio\\
\textbf{Title:} & DS$(r,1,1)$ for all $r \ge 2$: a direct proof via Thm 3.12 and Sub-Lemma Z\\
\textbf{Source:} & Corresponds to work-in-progress commit \texttt{e99cb59}\\
& {\footnotesize(\texttt{proofs/2026-09-26-day207-DS-r11-proved.md})}\\
\textbf{Registry:} & \texttt{hikita-star-dominance-support.json} nodes\\
& \texttt{ds-lambda-211-via-p2Y-pieri}, \texttt{ds-lambda-r11-all-r}.\\
& Trust: proved (self-proved, not peer-checked)\\
\textbf{Scripts:} & \texttt{scripts/day207/ds\_r11\_direct\_check.py} (+\,\texttt{.log}, ALL OK),\\
& \texttt{scripts/day207/ds\_r11\_route\_consistency.py} (+\,\texttt{.log}, ALL OK)\\
\textbf{Closes:} & \texttt{memory/questions/q-DS-211-promotion.md}
\end{tabular}
\end{center}

\medskip
\noindent In the registry, \texttt{ds-lambda-211-via-p2Y-pieri} is now proved and \texttt{ds-lambda-r11-all-r} is new.

\section{Statement}\label{sec:statement}

The conventions are those of \texttt{scripts/day198/p2Y\_er.py::build\_action}, which are the same as Day~205~\S1 and Day~206b~\S0.
\begin{itemize}
\item $s := q^{-1}$.
\item $\tn{n} := \tn{n}_t$, with $\tn{n}=0$ for $n\le 0$.
\item $e_\mu$ is the ordinary product of the $e_{\mu_i}(X_1,\dots,X_m)$.
\item $\st$ is Hikita's product (2503.23597, Def.~3.4). Define $C_r := e_1\st e_1\st e_r$.
\end{itemize}

\noindent\textbf{DS$(r,1,1)$} is the length-3 slice of DS-triangular (Day~196~\S2) at $\lambda=(r,1,1)$. It asserts two things:
\begin{itemize}
\item $C_r$ lies in the span of the $e_\mu$ with $\mu\succeq(r,1,1)$, $\mu\vdash r+2$. For $r\ge2$ these $\mu$ are exactly $(r+2)$, $(r+1,1)$, $(r,2)$ and $(r,1,1)$.
\item The coefficient of $e_{(r,1,1)}$ is $q^{-n((r,1,1))}=q^{-3}$, and every other coefficient vanishes at $q=1$.
\end{itemize}

\begin{theorem}\label{thm:main}
For every $r\ge1$ and $m\ge1$, the following identity holds in $\Lambda_m$:
\begin{multline*}
C_r = s^3\,e_{(r,1,1)} + s^2(1-s)\tn{2}\,e_{(r,2)} + s(1-s)\bigl(\tn{r+1}+t\tn{r}+s\bigr)\,e_{(r+1,1)}\\
 + (1-s)^2\tn{r+2}\bigl(\tn{r+1}+s\bigr)\,e_{(r+2)}.
\end{multline*}
\end{theorem}

Now take $r\ge2$ and $m\ge r+2$. The four $e_\mu$ are then distinct and linearly independent, so this is \emph{the} $e$-expansion. In particular:
\begin{itemize}
\item The support is exactly the DS-interval of $(r,1,1)$. All four coefficients are nonzero.
\item The leading coefficient is $s^3=q^{-3}$.
\item The three off-diagonal coefficients carry a factor $(1-s)$, so they vanish at $q=1$.
\end{itemize}
Hence \textbf{DS$(r,1,1)$ holds, in its sharpened triangular form, for all $r\ge2$.} It also holds in the stable limit.

At $r=2$ the formula reproduces the Day~196/198 table:
\[
q^{-3},\quad \frac{(q-1)(t+1)}{q^3},\quad \frac{(q-1)(2qt^2+2qt+q+1)}{q^3},\quad \frac{(q-1)^2(t+1)(t^2+1)(qt^2+qt+q+1)}{q^3}.
\]

\section{Proof (direct route)}\label{sec:proof}

\paragraph{Step 0: $\st$ is $e_1(Y)\bullet$.} This is Day~205 R0. Since $Y_i\bullet1=X_i$, we have $e_1(Y)\bullet1=e_1$ and $\qq^{-1}(e_1)=e_1(Y)$. For any $G$, $e_1\st G=(e_1(Y)\,\qq^{-1}G)\bullet1=e_1(Y)\bullet G$. Applying this twice:
\[
C_r = e_1\st(e_1\st e_r) = e_1(Y)\bullet\bigl(e_1(Y)\bullet e_r\bigr).
\]
This uses Hikita Def.~3.4 and the bijectivity of $\qq_{(m)}$ (cited; the source is at verified-quote). $\st$ is transported from the commutative product on $\Lambda(Y)$, so it is associative and commutative, and the order of the factors does not matter.

\paragraph{Step 1: Thm~3.12 in operator form, at $r$ and at $r+1$.} For all $n\ge1$:
\begin{equation}\label{eq:Tn}\tag{$T_n$}
e_1(Y)\bullet e_n = (1-s)\tn{n+1}\,e_{n+1} + s\,e_1e_n.
\end{equation}

\begin{proof}[Proof, from scratch]
Day~205 Step~A gives $e_1(Y)=\sigma_m\pi$ on symmetric $F$. Day~205 Step~B gives $\pi(e_n)=X_1e_n(\mathrm{tail})+sX_1^2e_{n-1}(\mathrm{tail})$. Day~205b~\S3 gives:
\begin{itemize}
\item $\sigma_m[X_1e_n(\mathrm{tail})]=M_1(n)/(1-t)=\tn{n+1}e_{n+1}$;
\item $\sigma_m[X_1^2e_{n-1}(\mathrm{tail})]=M_2(n-1)/(1-t)=-\tn{n+1}e_{n+1}+e_1e_n$.
\end{itemize}
Adding these with weights $1$ and $s$ gives \eqref{eq:Tn}.
\end{proof}

This is Hikita's Thm~3.12, re-derived. It is also the ``anchor'' check in Day~205 \texttt{reduction\_0}~(e).

Applying $(T_r)$:
\begin{equation}\label{eq:one}
C_r = (1-s)\tn{r+1}\cdot e_1(Y)\bullet e_{r+1} + s\cdot e_1(Y)\bullet(e_1e_r).
\end{equation}

\paragraph{Step 2: the two pieces.}
\begin{itemize}
\item $(T_{r+1})$ gives $e_1(Y)\bullet e_{r+1}=(1-s)\tn{r+2}e_{r+2}+s\,e_{(r+1,1)}$.
\item Sub-Lemma~Z (operator form) is proved in Day~205 Steps~A--D together with Day~205b (L1)--(L4). It holds for all $r\ge1$ and $m\ge1$:
\[
e_1(Y)\bullet(e_1e_r) = (1-s)^2\tn{r+2}e_{r+2} + (1-s)(t\tn{r}+s)e_{(r+1,1)} + s(1-s)\tn{2}e_{(r,2)} + s^2e_{(r,1,1)}.
\]
(Day~205b writes this as $(q-1)(qt\tn{r}+1)/q^2$ for the $(r+1,1)$ coefficient, which is the same thing.)
\end{itemize}

\paragraph{Step 3: assembly.} Substitute both pieces into \eqref{eq:one} and collect terms:
\begin{center}
\renewcommand{\arraystretch}{1.3}
\begin{tabular}{@{}llll@{}}
\toprule
$\mu$ & from $(1-s)\tn{r+1}\cdot(T_{r+1})$ & from $s\cdot$Z & total\\
\midrule
$(r+2)$ & $(1-s)^2\tn{r+1}\tn{r+2}$ & $s(1-s)^2\tn{r+2}$ & $(1-s)^2\tn{r+2}(\tn{r+1}+s)$\\
$(r+1,1)$ & $s(1-s)\tn{r+1}$ & $s(1-s)(t\tn{r}+s)$ & $s(1-s)(\tn{r+1}+t\tn{r}+s)$\\
$(r,2)$ & $0$ & $s^2(1-s)\tn{2}$ & $s^2(1-s)\tn{2}$\\
$(r,1,1)$ & $0$ & $s^3$ & $s^3$\\
\bottomrule
\end{tabular}
\end{center}
This proves the identity for every $m$.

\paragraph{Uniqueness.} Take $r\ge2$ and $m\ge r+2$. Then $(r+2)$, $(r+1,1)$, $(r,2)$ and $(r,1,1)$ are distinct partitions with largest part $\le m$. Since $e_1,\dots,e_m$ are algebraically independent (Macdonald I~(2.4)), the $e_\mu$ are linearly independent.

\paragraph{Nonvanishing.} At $q=2$ and $t=1$, the four coefficients are $1/8$, $1/4$, $\tfrac14(r+1+r+\tfrac12)$ and $\tfrac14(r+2)(r+\tfrac32)$. All are nonzero. \qed

\subsection*{Why this is not the Day~195/198 tautology}
The Day~198 note (\S6) found that iterating Thm~3.12 collapses to $C=C$. The reason was that $A=e_1\st(e_1e_r)$ was an unknown. Sub-Lemma~Z computes $A$ independently, from the level-1 action, and that is what breaks the circle. The route uses only $e_1(Y)$. It does not use $p_2(Y)$, Newton, $W_r$ or $\tau_r$.

\section{Second proof (Day~198 route, now also closed)}\label{sec:route2}

Newton's identity $e_1(Y)^2=p_2(Y)+2e_2(Y)$ holds in the commutative ring $\Lambda(Y)$. Acting on $e_r$ gives:
\[
C_r = p_2(Y)\bullet e_r + 2\,e_2(Y)\bullet e_r = p_2(Y)\bullet e_r + 2t\,W_r.
\]
The inputs are:
\begin{itemize}
\item Lemma~1 ($p_2(Y)$-Pieri, including $\tau_r$), proved as Day~206b~\S6;
\item $W_r=t^{-1}e_2(Y)\bullet e_r$, proved in Day~206b.
\end{itemize}
Support and leading coefficient then follow as in Day~198~\S5:
\begin{itemize}
\item $p_2(Y)\bullet e_r$ is supported on the DS-interval, and its $(r,1,1)$-coefficient is $q^{-3}$;
\item $W_r$ lives on the two-row interval, so its $(r,1,1)$-coefficient is $0$.
\end{itemize}
\texttt{ds\_r11\_route\_consistency.py} checks both routes against the closed form symbolically in $u=t^r$, which covers all $r$ at once. All four coefficients agree ($\Delta=0$).

Route~1 is preferred because it has fewer inputs. The only input it shares with Route~2 is the Day~205/205b machinery.

On the question file: the $D$ in Day~198's $(\st)$ is $D=e_2\st e_2=W_2$. It is not an ``$e_1\st(e_1\st e_2)$-type'' term. Sub-Lemma~Z enters Route~1 directly, and enters Route~2 only through Day~206b.

\section{Inputs and grades}\label{sec:grades}

{\small
\begin{center}
\renewcommand{\arraystretch}{1.25}
\begin{tabular}{@{}p{4.6cm}p{1.3cm}p{5.2cm}p{3.0cm}@{}}
\toprule
Input & Used in & Grade & Where\\
\midrule
R0: $e_1\st G=e_1(Y)\bullet G$ (Hikita Def.~3.4, $\qq$ bijective) & Route 1 & proved, modulo citation (2503.23597 at verified-quote) & Day 205 \S2 R0\\
Step A ($e_1(Y)=\sigma_m\pi$ on symmetric $F$), Step B ($\pi$-split) & Route 1 & proved & Day 205 \S2\\
Lemma 1 (closed form of $\sigma_m$) and Lemma 2 ($(1-t)S_n=q_n$); $M_1$, $M_2$ & Route 1 & proved & Day 205b \S\S1--3\\
Thm 3.12, operator form $(T_n)$ & Route 1 & proved (re-derived here from the three rows above); also Hikita, verified-quote & \S2 Step 1\\
Sub-Lemma Z, operator form & Route 1 & proved (Clio UID 286 independently proved (L1)--(L4)) & Day 205, Day 205b\\
Newton decomposition & Route 2 & proved & Day 198 \S3\\
Lemma 1 $=p_2(Y)$-Pieri, including $\tau_r$ & Route 2 & proved & Day 206b \S6\\
$W_r=e_2\st e_r$ & Route 2 & proved & Day 206b\\
\bottomrule
\end{tabular}
\end{center}
}

The weakest link is the same Hikita interface (R0) that every node in this tree shares. It concerns the convention match between the code's $Y_i$ and Hikita's eq.~(3), and the bijectivity of $\qq$. As a statement about the implemented operator $e_1(Y)^2$, the theorem is fully proved.

\section{Verification}\label{sec:verification}

The script \texttt{scripts/day207/ds\_r11\_direct\_check.py} uses the flint AHA engine from\\ \texttt{day205/k3\_fast\_pipeline.py}. Its conventions are verbatim from Day~198. For each $r=2,\dots,5$ and $m\in\{r+2,r+3\}$, it computes three quantities directly and compares their $e$-expansions exactly with the claimed formulas:
\begin{itemize}
\item $e_1(Y)\bullet e_{r+1}$;
\item $e_1(Y)\bullet(e_1e_r)$;
\item $C_r=e_1(Y)\bullet(e_1(Y)\bullet e_r)$.
\end{itemize}
Result:
\begin{itemize}
\item 24/24 checks pass (Thm~3.12, Z and $C_r$ at each of the 8 $(r,m)$ pairs).
\item The support is always $\{(r+2),(r+1,1),(r,2),(r,1,1)\}$.
\item The leading coefficient is always $q^{-3}$.
\item Log: \texttt{ds\_r11\_direct\_check.log}.
\end{itemize}
\texttt{ds\_r11\_route\_consistency.py} checks that Route~1 and Route~2 both equal the closed form, symbolically in $r$.

\section{What this does not do}\label{sec:scope}

\begin{itemize}
\item It does not touch DS at $(2,2,1)$, $(2,1,1,1)$, $(1^4)$, or at length $\ge3$ in general. Those remain \texttt{computed} (Day~196).
\item The natural next case is $(r,2,1)=e_1\st e_2\st e_r$. It is $e_1(Y)\bullet W_r$ (with $W_r=e_2\st e_r$ from Day~206b), which needs $e_1(Y)\bullet(e_1e_{r+1})$ (Sub-Lemma~Z at $r+1$), $e_1(Y)\bullet(e_2e_r)$ and $e_1(Y)\bullet e_{r+2}$. So the only missing piece is a ``Sub-Lemma~Z$_2$'' for $e_1(Y)\bullet(e_2e_r)$. That is the same $\pi$-split and $\sigma_m$ machinery with $g=e_{r-1}(\mathrm{tail})$ and $e_2(\mathrm{tail})=e_2-X_1e_1+X_1^2$. Cheap next target.
\end{itemize}

\end{document}
